\section{Freight and Transport}

\begin{frame}{Foreign Trade: Freight}
    \mycomment{
    International trade begins with someone sensing an opportunity, and then an importer and an exporter teaming up together. A business relation is established, and a contract is concluded. This contract must include quite a lot of fine print:

    \begin{itemize}
        \item what exactly is traded?
        \item in which quantity, price, and when?
        \item who pays when?
        \item who transports the wares?
        \item how is the ware insured?
        \item ...
    \end{itemize}
    
    % We will be looking at exactly how to transport things from A to B, and look at some of the jargon in shipping, including road and rail. 
    
    % Typically, our delivery will get from the origin in a mine, farm or factory (upstream) to a port, then be shipped (as freight), where it crosses a country border, over international waters into the realm of the importing country, where it crosses another border and may become subject to tariffs or other types of import restrictions, and be forwarded via dry ports, road or rail to the refinery, factory, warehouse, and finally end consumer.

    % A financial flow happens in the opposite direction; this is the direction in which the order is placed. There are, however intermediates. It needs to be clarified who pays for the shipping (f.o.b. or c.i.f.), and when and where property is transferred.
    }
    \showmytikz[0]{../shared/tikz/tikz-freight.tex}
    
\end{frame}


\begin{frame}{History of Trade}
    \mytext{
    \begin{itemize}
        \item Bronze age (Uluburun shipwreck 1400 v.Chr.)
        \item Roman Empire, Umayyad and Abbasid Caliphates
        \item Hanseatic League, Silk Road, Venice and Genoa
        \item Pre-industrial Colonialism: Portugal, Spain
        \item Industrial Colonialism: Britain
        \item World Wars; / Bretton Woods System, GATT
        \item Cold War; Third World; developing countries
        \item Oil Crisis, floating currencies
        \item WTO, Globalisation and new imbalances
        \item Recent developments
    \end{itemize}
    }
    \mycomment{
        Progress reduced transportation costs; the discovery of America, the Suez canal, the steam boat. 

        Along came innovations like paper money, the banking system, stock exchanges, exchange traded liablities, insurances etc.

        States and violence played a role: securing trade routes, establishing trading posts, monopolisation of trade, slavery and coercion, etc.
    
    Literature tipps: \cite{polanyi44}, \cite{beckertcapital}
    }

\end{frame}

\begin{frame}{Where Do We Trade?}
\mycomment{ 
    `Over the counter' refers to a normal, bilateral trade: a seller and a buyer make a private contract. 

    In contrast, many commodities are traded on exchanges. While I can buy a ton of coffee over the counter, I am likely to get the best price and have the smoothest operation if I buy a standardised contract on a mercantile exchange.

    Contract on mercantile exchanges have standardised physical and financial fulfilment. Often, they come with third party guarantees. The structure, ownership etc. of these exchanges vary widely.
    }
\mytext{
    \begin{itemize}
        \item Over the Counter (OTC)
        \item Stock Exchanges / Commodity Exchanges
            \begin{itemize}
                \item CME (Chicago, commodities)
                \item ICE (International Contental Exchange), LCE (London Commodity Exchange)
                \item London Metal Exchange
                \item Euronext (multiple places)
                \item New York Stock Exchange \& Nasdaq
                \item Japan Exchange Group (Tokyo Stock Exchange)
                \item EEX (Leipzig, energy)
            \end{itemize}
    \end{itemize}
}
\end{frame}



\begin{frame}{Goods / Commodities}
\mycomment{
We typically distinguish between raw materials, commodities, and merchandise, where each category is a sub-category of the other. So the wares imported and exported can have different degree of manufacture, raging from freshly mined or harvested produce to sophisticated machinery.
}
\mytext{
    \mylonglist{
    
    Commodities
    \begin{itemize}
        \item `Hard' 
        \begin{itemize}
            \item Industrial Metals
            \item Precious Metals
            \item Energy
        \end{itemize}
        \item 'Soft' 
        \begin{itemize}
            \item Corn, rice, etc.
            \item Animal products
        \end{itemize}
    \end{itemize}
    Goods / Merchandise
    \begin{itemize}
        \item Manufactured / consumer goods
        \item Machinery
        \item Arms, automotives
        \item Electronics
        \item ...
    \end{itemize}
    }
}
\end{frame}

\begin{frame}{Value Chain, Upstream, Downstream}
    \showmytikz{../shared/tikz/tikz-valuechain.tex}
    \mycomment{
    Raw material is mined (or grown), then refined. Refined commodities are then used to produce intermediate goods, which are combined to manufactured wares. These are then branded and distributed.

    At all stages, costs occur, and value is added. Production requires not only labour, but capital in the form of machinery etc. Distribution is very labour intensive, and includes sales and marketing.    
    
    In most cases, production and distribution contribute the largest part to the value of the final consumer good. 
    }
\end{frame}


% \begin{frame}{Special Case: Energy}
%     \mytext{
%     \begin{itemize}
%         \item Actual energy only in the form of electricity
%         \item kinetic and potential energy
%         \item chemical energy is potential energy
%         \item Electricity (grid based)
%         \item Coal, oil and gas: pipelines vs. cargo
%     \end{itemize}
%     }
% \end{frame}

\begin{frame}{Other Tradables}
\mycomment{
Any private good can be traded, in principle. We will focus on material goods - and the trade balance - in this class. It is, however, worth keeping in mind that the financial and capital accounts, that the import and export of credit and liability is technically also import and export. Any physical trade is mirrored by a financial transaction in the opposite direction. And while physical goods are often subject to tariffs, capital transfers are (today) typically not taxed.
}
\mytext{
    \mylonglist{
    
    \begin{itemize}
        \item Financial Services
        \item IT services
        \item Legal services
        \item Usage rights, licenses
        \item Data transmission / satellite
        \item Internet based
        \item Big data?
        \item Energy
        \item Currency
        \item Financial Instruments
        \item ...
    \end{itemize}
    }
}
\end{frame}

 
 

\begin{frame}{Fulfilment and Transfer of Property}
    \mytext{
    Typical (trading) contracts have two sides:
    \begin{itemize}
        \item physical fulfilment, that is: the seller gives the buyer the traded good
        \item financial fulfilment, that is: the buyer pays the seller
    \end{itemize}
    The time and place for both need to be specified.
    \begin{itemize}
        \item FOB free on board
        \item CIF cost insurance freight
        \item Factory door / doorstep
    \end{itemize}
    }
    \mycomment{
    The actual moment in which ownership, and hence responsibility, is transferred, is one of the most important aspects of any contract. Are the wares on board of a ship still in the seller's possession, or already property of the buyer? Is the cotton in the factory, where it is spun into yarn, property of the yarn manufacturer, or is the spinner delivering a service to the yarn owner?
    }
    
\end{frame}


\begin{frame}{Incoterms}
    \mycomment{Incoterms are definitions by the International Chamber of Commerce. They typically consist of three-letter abbreviation common in international trade, specifying who of the contracting partners in in charge of which part of the delivery chain.    
    }
    \mytable{
    \begin{tabular}{ll}
        \toprule
         EXW & ex works  \\
         FCA & free carrier \\
         FAS & free alongside ship \\
         \textbf{FOB} & \textbf{free on board} \\
         \midrule
         CFR & cost and freight \\
         \midrule
         \textbf{CIF} & \textbf{cost insurance freight} \\
         DAP & delivered at place \\
         DPU & delivered at place unloaded \\
         CPT & carriage paid to \\
         CIP & carriage insurance paid \\
         DDP & delivered duty paid \\
         \bottomrule
    \end{tabular}
    
    }
    
\end{frame}

\begin{frame}{Purchase Order Lead Time (POLT)}
    \mycomment{
    Trade does not happen without time lags, and is not without risk. The time lags require careful planning, which also include the insurance against risk. Among these risks is the change of price (and that of currency) during the voyage. To insure agains this, futures and forward contracts have been invented.

    `Buy now, receive tomorrow' should
    \begin{itemize} 
        \item normally be more expenseive than spot
        \item due to cost of storage etc.
        \item if so: `contango'
        \item if not: `backwardation'
    \end{itemize}

    \textit{Think about you buying my car today but picking it up only in a month's time.}
    }
    \mytext{
    
    \begin{itemize}
        \item Placing order
        \item Confirm order (backorder?)
        \item Ship date (2-day shipping?)
        \item Estimated shipping date (ETA)
        \item Delivery at port / customs
        \item Delivery at factory
        \item Fulfilment 
        \item Delivery exceptions?
        \item Return date?
    \end{itemize}
    }

\end{frame}

\begin{frame}{Insurances and Futures}
    \mycomment{
    \begin{itemize}
        \item Invoice creation date ($\neq$ billing date)
        \item Delivery date 
        \item Payment?
        \item Order tracking?
        \item Insurance?
    \end{itemize}
    What could cause delays or go wrong?
    \begin{itemize}
        \item Holidays
        \item Weather-impacted slowdowns
        \item Customs (can take days or weeks)
        \item Currency revaluation
        \item Price movements
        \item Force majeure
    \end{itemize}
    }
    \mytext{
    
    \begin{itemize}
        \item Forward (OTC)
        \item Futures (publicly traded)
        \item Put and Call options
        \item Insurance
        \item ...
        \item Counterparty risk
    \end{itemize}
    }
 
    
\end{frame}

\begin{frame}{Export Credit Agencies (ECA)}
\mycomment{
    Many countries support their exporters with credits and insurances. This is particularly true for the richer countries, that often have governmental credit agencies in place covering the often significant risks involved in the export of complex machinery, but also of FDI.
}
\mytext{
    \begin{itemize}
        \item export insurance
        \item Germany: \textit{Hermes-Deckung, KfW}
        \item cf. \textit{ungebundene Finanzkredite, UFK}
        \item Janan: \textit{JICA}
        \item inofficial OECD agreement
        \item for exports of goods and capital (FDI)
        \item Level playing field?
        \item Good for developing countries?
        \item `Country risk' (Peren-Clement index)
    \end{itemize}
}
\end{frame}


\begin{frame}{International Standards and Arbitration}
\mycomment{
    National and international trade associations provide standardised contracts and arbitration / dispute settlement. An important framework is called "Incoterms".
    
    Out-of-court-settlements are often preferable to settlement in courts. Thus, many contracts include provisions about dispute settlement out of court, specifying the procedure and organisations involved beforehand.
    }
\mytext{
    \begin{itemize}
        \item United Nations Convention on Contracts for the International Sale of Goods (CISG) \href{https://uncitral.un.org/sites/uncitral.un.org/files/media-documents/uncitral/en/a-conf-97-19-ocred-eng.pdf}{(external link)}
        \item International Chamber of Commerce ICC  \href{https://iccwbo.org/}{https://iccwbo.org/}
        \item German Chamber of Commerce and Industry \href{https://www.dihk.de/}{https://www.dihk.de/} (Every spring (May): World Business Outlook)
        

    \end{itemize}
}
\end{frame}


\begin{frame}{Quality Assurance - Samples and Prototypes}
\mytext{
    The CE sign is required to bring goods into circulation within the EU. By printing it on the product the buyer assures that the product complies with all relevant EU regulations.
    
    \showmybitmap{../shared/bitmaps/ce.png}
}
\mycomment{
    \begin{itemize}
        \item initial sample
        \item prototypes?
        \item inspection?
        \item CE sign - self-commitment
    \end{itemize}
}
\end{frame}


%%%%%%%%%%%%%% freight


\begin{frame}{Road and Rail Transport}
    \mycomment{Once the contract is in place, fulfilment begins. That means that the wares are shipped. The first leg is from the factory or warehouse to a port; it is similar to the last leg, from the port to the importer's warehouse or factory.}
    \showmytikz[1,2]{../shared/tikz/tikz-freight}
    
\end{frame}


\begin{frame}{Packaging Industry}
    \mycomment{The packaging industry is a large industry with close ties to the logistics industry, which often works with re-packaging. Dry ports often serve as re-packaging hubs (intermodal hubs), where e.g. merchandise is unloaded from containers and re-packaged on Euro-pallets. The (chinese) pallets are discarded, the goods re-wrapped in clingfoil. The amount of waste is significant, though to a good deal easy to recycle.}
    \mytext{
    \begin{itemize}
        \item cling foil, pallets
        \item postal packaging, card bord boxes
        \item filling material (plastic, paper, corn)
        \item bubble wrap
        \item bottling, barrels
        \item containers
        \item ...
    \end{itemize}
    }
\end{frame}




\begin{frame}{Global Logistics / Cargo}
\mycomment{On top of global logistics, there is of course the logistics of individual companies, the organisation of which is part of business administration. Here are some further keywords to consider or to know when dealing in international trade.}
\mytext{
    \begin{itemize}
        \item Warehouse management?
        \item Just in time?
        \item Third-party logistics, 4PL 
        \item freight forwarder
        % \item FoB (free on board)
        % \item CIF (costs insurance freight)
        \item Logistics Performance Index
        \item Parcel / postal delivery
        \item LCL (less than container load)
        \item EXW ex works (factory door)
    \end{itemize}
}
\end{frame}


%%% ROAD


\begin{frame}{Road}
\mytext{
    \begin{itemize}
        \item Containers,  Tankers
        \item Less-than-truckload (LTL) shipinng
        \item Semi-trailer (`Sattelauflieger')
        \item customs transit system:
        \item IRU; TIR system / carnet TIR  (Transports Internationaux Routiers)
    \end{itemize}
    \showmybitmap{../shared/bitmaps/Containerchassis_40fuss.jpg}
    }
    \mycomment{
    \includegraphics[width=80pt]{../shared/bitmaps/Tir_Plate.svg.png}

    TIR plate 
    }
\end{frame}

\begin{frame}{Rail and EURO Pallet}
    \mytext{
    UNECE: Unified Railway Law (UN)

    EU railways, since 1968: 
    European Pallet Pool (EPP) swaps pallet for pallet

    EPAL European Pallet Association, reburse-system

        \showmybitmap{../shared/bitmaps/Europalette.png}
    }
    \mycomment{
    The metric sizes of EPAL are obviously incompatible with the imperial sizes used in container shipping.

    \begin{itemize}
        \item 800mm x 1200mm x 144mm
        \item app. 25kg
    \end{itemize}

    International freight typically uses single-use pallets. Consequentially, a lot of re-packaging has to be done.

    \begin{itemize}
        \item Spare space in containers
        \item One-way pallets
        \item compatible intermodal container
    \end{itemize}    
    }
\end{frame}



\begin{frame}{Euro Container}
    \mytext{
    German Association of Automotive Industry (VDA)
    \begin{itemize}
        \item KLT Box, Eurobox
        \item Fits EPAL
        \item $600mm\times 400mm$ 
        \item $400mm\times 300mm$        
    \end{itemize}
    \showmybitmap{../shared/bitmaps/eurocontainer.jpg}
    }
    \mycomment{These come in a variety of colours, like red for meat, basket shape for bakery products, with lids or without lids, open handles, closed handles, etc.

    Smaller ones stack on larger ones, but larger ones don't stack on smaller ones. On paletts they are typically held together by cling foil.
    }
\end{frame}


\begin{frame}{Flexitank: Bulk Liquid Food (and Petrol) transport}
    \mycomment{Liquids can be transported in large flexible tanks. However, juice. soft drinks and similar substances are transported without in powder form the water content. Then, the water content remaining is a crucial quality to compare quantities.

    }
    \mytext{
        
    \begin{itemize}
        \item Made out of rubber
        \item Tank fits inside standard container
        \item Tank can be fitted onto rail and road
        \item intermodel freight transport
        \item (often) food-grade (e.g. wine, cooking oil)
    \end{itemize}
    }
\end{frame}

\begin{frame}{Barrels (Oil Barrels)}
    \mycomment{Some goods are shipped in quantities that use traditional, non-metric units, most prominently crude oil. However, the quantity measure "barrel" is common, even though petrochemicals may not be packaged at all.
    }
    \mytext{
    \begin{itemize}
        \item traditional barrels, wine barrels
        \item oil barrel: 42 US gallons, $=158.99\text{l}$ 
        \item unit is $b$ or $bbl$
        \item at $60^oF = 15.6^oC$ and $14.696 psi$
    \end{itemize}
    \showmybitmap{../shared/bitmaps/oilbarrel.png}
    }
\end{frame}

\begin{frame}{Pipelines}
    \mycomment{
        A common land (and to some extent, sea bed) based means of transport for liquids are pipelines. These can be owned by a single operator. However, at least in Europe, pipelines are infrastructure similar to electricity grids: multiple traders can use the same pipeline to deliver oil and gas. The substance transported is always the same; however, ownership, price etc. can be different.
    }
    \mytext{
    \begin{itemize}
        \item Transport of Liquids
        \item by far safest and cheapest
        \item critical infrastructure
        \item Crude Oil
        \item Use as transport network: no monopoly (EU)
    \end{itemize}
    }
\end{frame}

%%% SHIPS 


\begin{frame}{Shipping}
    \mycomment{
    The vast majority of global trade takes place via shipping. Special care must be given for who pays for the shipping, and who insures the wares during shipping. 
    }
    \showmytikz[1,2,3]{../shared/tikz/tikz-freight}    
\end{frame}


\begin{frame}{Cargo Ships}
    \mytext{
    \begin{itemize}
        \item breakbulk cargo (`Stückgut')
        \item bulk cargo (`Schüttgut')
        \item containerised cargo
        \item bulk carrier, tank ship
    \end{itemize}
    % \includegraphics[height=0.4\textheight]{}
    \showmybitmap{../shared/bitmaps/Sabrina_I_cropped.jpg}
    }
\end{frame}


\begin{frame}{A Bit of History}
    \mycomment{People have been trading for a very long time. The advances in navigation, the discovery of America and the colonisation of much of the world by Europeans have significantly increased global trade.}
    \mytext{
    \begin{itemize}
        \item gangs, longshoremen, dock workers
        \item dry bulk cargo, break bulk cargo
        \item Malcom McLean: containerisation
        \item Decline in workers needed
        \item Increase in cargo
        \item Decline of inland ports
        \item Increase of `dry ports' (road cargo)
    \end{itemize}
    }
\end{frame}


\begin{frame}{ISO Container}
    \mytext{Inner measures:
        \showmytikz{../shared/tikz/tikz-iso668container}
    }
    
    \mycomment{
    ISO Container have only been around since about the 1960ies. Containerisation of global trade has lead to a massive increase in the trade of merchandise.
    
    \begin{itemize}
        \item 8 feet $\times$ wide $=2.44m$ (max. ca. $30t$)
        \item 8 feet 6 inches high $=2.49m$
        \item 20 or 40 feet long $6.10 or 12.19m$
        \item 9 feet 6 inches high $=2,90m$ (High Cube HC or HQ)
    \end{itemize}
    % \includegraphics[height=0.5\textheight]{images/1024px-Container_Augsburg.jpg}
    }
    
\end{frame}

\begin{frame}{TEU, Ship Sizes}
    \mycomment{There are a number of noteworthy units and abbreviations regardings commercial ships.}
    \mytext{
    TEU (Twenty Foot Equivalent Units)
    \begin{itemize}
        \item FCL full container load
        \item LCL less than container load: \textit{groupage}
    \end{itemize}
    DWT deadweight tonnage
    \begin{itemize}
        \item liner
        \item tramp
    \end{itemize}

    }
    
\end{frame}

\begin{frame}{Cargo Ships}
    \mytext{HMM Algericas, one of the 12 largest container vessels.
    
    \showmybitmap{../shared/bitmaps/hmm_algericas_Dietmar_Hasenpusch.jpg}
    }
    \mycomment{
    IMO 9863297\\
    length: 399.9m\\
    width: 61.03m\\
    GT 228.283\\
    DWT 232.606\\
    23.964 TEU\\
    Operator: HMM
    }
\end{frame}

\begin{frame}{Size limits}
\mycomment{Sizes of ships are often given by the maximum sizes of narrow passages they can transverse. The most important are the Panama Canal (1904), the Suez Canal (1869), St. Lawrence Seaway (Great Lakes, Seawaymax), the Strait of Malacca, and the size of Chinese ports (Chinamax). Everythign larger goes around Cape Agulhas / Cape Horn.}
\mytext{
    \begin{itemize}
        \item Panamax
        \item Suezmax
        \item Seawaymax
        \item Malaccamax
        \item Chinamax
        \item Capesize
    \end{itemize}
}
\end{frame}


\begin{frame}{LNG / Pressure Tanks}
    \mytext{
    LNG is liquified natural gas, hence (98\%) Methane. Do not confuse LNG with LPG (liquified petroleum gas) which is mostly propane/butane gas, aka camping gas. The latter is used for cooking in many poorer countries.
    \begin{itemize}
        \item LNG (liquified natural gas) liquidity:
        \item high pressure
        \item low temperature ($-165$°)
        \item typically spherical
    \end{itemize}
    \showmybitmap{../shared/bitmaps/Shahamah_LNG_carrier.jpg}
    }
    \mycomment{LNG becomes liquid under very low temperatures and high pressure, hence the spherical containers. The cooling on the transport consumes part of the stored energy.}
\end{frame}

%%% AIR FREIGHT

\begin{frame}{Air Freight and Passenger Transport}
    \mycomment{Aeroplanes use specialised containers, and these may not be interoperable between aircraft types. Boeing and Airbus have different systems.
    
    Air traffic is dominated by passenger traffic, whereas freight is a supplement and somewhat a by-product, particularly since the advent of electronic mail. Within international passenger traffic, air traffic dominates all other modes of transport. As such, air traffic is often seen as part of the tourism industry, which is a significant sector in many economies. 

    In comparison to tourism, the number of people (workers) re-locating (migration) is negligible. However, once migrated, workers remittances to their home country are sometimes economically significant for the receiving (poorer) economies.
    }
    \mytext{
    \begin{itemize}
        \item Unit load device (ULD), ca. $11.52m^3$
        \item ULD pallet (alumininum)
        \item ULD container (can or pod)
    \end{itemize}
    \showmybitmap{../shared/bitmaps/ULD-1.png}
    }

    % \includegraphics[height=0.5\textheight]{images/1280px-Unloading_JAL_747.jpg}
    
\end{frame}




\begin{frame}{Container Port Sizes}
\mytable{
        \begin{tabular}{rllrl}
             \toprule
             Nr. & Port & Country          & MTEU & \\
             \midrule 
             1 & Shanghai & China          & 51  & \myblock[0.5ex]{51}\\
             2 & Singapore & Singapore     & 41  & \myblock[0.5ex]{41}\\
             3 & Ningbo-Zhoushan & China   & 39  & \myblock[0.5ex]{39}\\
             4 & Shenzen & China           & 33  & \myblock[0.5ex]{33}\\
             5 & Quingdao & China          & 30  & \myblock[0.5ex]{30}\\
             6 & Guangzhou & China         & 26  & \myblock[0.5ex]{26}\\
             7 & Busan & South Korea       & 24  & \myblock[0.5ex]{24}\\
             8 & Tianjin & China           & 23  & \myblock[0.5ex]{23}\\
             9 & Dubai & UAE               & 16  & \myblock[0.5ex]{16}\\
             10 & Port Klang & Malaysia    & 15  & \myblock[0.5ex]{15}\\
             11 & Rotterdam  & Netherlands & 14  & \myblock[0.5ex]{14}\\
             12 & Hong Kong  & China       & 14  & \myblock[0.5ex]{14}\\
             13 & Antwerp-Bruges & Belgium & 14  & \myblock[0.5ex]{14}\\
             \bottomrule
        \end{tabular}
        \par \cite{lloydslist}
    }
\mycomment{} %% empty, but so that the reference gets produced.
\end{frame}

\begin{frame}{Container Port Sizes}
\mycomment{}
\mytable{
        \begin{tabular}{rllrl}
             \toprule
             Nr. & Port & Country & MTEU 2025 & \\
             \midrule
             13 & Antwerp-Bruges & Belgium   & 14 & \myblock[0.5ex]{14}\\
             14 & Xiamen & China             & 12 & \myblock[0.5ex]{12}\\
             15 & Tanjung Pelepas & Malaysia & 12 & \myblock[0.5ex]{12}\\             
             16 & Los Angeles & USA          & 10 & \myblock[0.5ex]{10}\\
             17 & Tanger Med & Morocco       & 10 & \myblock[0.5ex]{10}\\
             18 & Suzhou & China             & 10 & \myblock[0.5ex]{10}\\
             19 & Long Beach & USA           & 10 & \myblock[0.5ex]{10}\\
             20 & Laem Chabang & Thailand    & 10 & \myblock[0.5ex]{10}\\

             \bottomrule
        \end{tabular}
        \par Source: 2025 Lloyd's List Intelligence \cite{lloydslist}
        }

\end{frame}
\begin{frame}{Container Port Sizes}
\mycomment{}
\mytable{
        \begin{tabular}{rllrl}
             \toprule
             Nr. & Port & Country & MTEU 2025 & \\
             \midrule
             24 & New York/New Jersey & USA         & 8 & \myblock[0.5ex]{8}\\
             25 & Mundra & India                    & 8 & \myblock[0.5ex]{8}\\
             26 & Tanjung Priok & Indonesia         & 8 & \myblock[0.5ex]{8}\\
             27 & Hamburg & Germany                 & 8 & \myblock[0.5ex]{8}\\
             28 & Colombo & Sri Lanka               & 8 & \myblock[0.5ex]{8}\\
             29 & Hai Phong & Vietnam               & 7 & \myblock[0.5ex]{7}\\
             30 & Cai Mep & Vietnam                 & 7 & \myblock[0.5ex]{7}\\
             31 & Jawaharlal Nehru (Mumbai) & India & 7 & \myblock[0.5ex]{7}\\
             \bottomrule
        \end{tabular}
        \par \cite{lloydslist}
        }
\end{frame}

\begin{frame}{Container Port Sizes}
\mycomment{}
\mytable{
        \begin{tabular}{rllrl }
             \toprule
             Nr. & Port & Country & MTEU 2025 & \\
             \midrule
             38 & Valencia & Spain             & 5 & \myblock[0.5ex]{5}\\
             39 & Abu Dhabi & UAE              & 5 & \myblock[0.5ex]{5}\\
             \midrule
             43 & Piraeus & Greece             & 5 & \myblock[0.5ex]{5}\\
             44 & Algericas & Spain            & 5 & \myblock[0.5ex]{5}\\
             45 & Tokyo & Japan                & 5 & \myblock[0.5ex]{5}\\
             46 & Bremen/Bremerhaven & Germany & 4 & \myblock[0.5ex]{4}\\
             70 & Yokohama & Japan             & 3 & \myblock[0.5ex]{3}\\
             76 & Nagoya & Japan               & 3 & \myblock[0.5ex]{3}\\
             \bottomrule
        \end{tabular}
        \par \cite{lloydslist}
        }
\end{frame}

\begin{frame}{Largest Container Shipping Companies}
\mycomment{~}
    \mytext{~}    
    \mytable{
    \begin{tabular}{llrrrl}
        \toprule
        Company & Headquarters & MTEU &  Vessels & Share & \\
        \midrule
        MSC & Geneva & 6.4 & 889          & 20.2\% & \myblock[0.6ex]{20.2}\\
        Maersk & Copenhagen & 4.5 & 735   & 14.3\% & \myblock[0.6ex]{14.3}\\
        CMA CGM & Marseille & 3.8 & 664   & 12.7\% & \myblock[0.6ex]{12.7}\\
        COSCO & Shanghai & 3.4 & 519      & 10.6\% & \myblock[0.6ex]{10.6}\\
        Hapag-Lloyd & Hamburg & 2.3 & 300 & 7.4\% & \myblock[0.6ex]{7.4}\\
        ONE & Tokyo/Sing. & 2.0 & 256     & 6.2\% & \myblock[0.6ex]{6.2}\\
        Evergreen & Taoyuan & 1.8 & 225   & 5.7\% & \myblock[0.6ex]{5.7}\\
        HMM & Seoul & 0.9 & 83            & 2.9\% & \myblock[0.6ex]{2.9}\\
        ZIM & Haifa & 0.8 & 130           & 2.5\% & \myblock[0.6ex]{2.5}\\
        Yang Ming & Keelung & 0.7 & 98    & 2.2\% & \myblock[0.6ex]{2.2}\\        
        \bottomrule
    \end{tabular}
    }

\end{frame}

\begin{frame}{Largest Oil Tanker Companies}
\mycomment{The size of the crude oil transport market is difficult to estimatate, since many tankers are sailing under convenience flags and ownership is often intransparent. Tanker content can be loaded and off-loaded from at sea, be it from offshore drilling platforms or from one tanker to the other. This is sometimes necessary as the largest vessels cannot moor in ports.}
    \begin{itemize}
        \item Mitsui OSK Lines
        \item Teekay Corp
        \item Euronav NV
        \item Scorpio Tankers Inc
        \item Frontline Ltd
        \item DTH Holdings
        \item Tsakos Energy Navigation
        \item ...
    \end{itemize}
\end{frame}

\begin{frame}{Vessel Finder}
    \mycomment{There are online tools to find the exact position of any registered vessel, which is of course interesting - not only if you want to monitor where your grand parents are on their cruse ship.}
    \showmybitmap[https://www.vesselfinder.com/]{../shared/bitmaps/vesselfinder.png}

\end{frame}



\begin{frame}{Law of the Seas}
    \mycomment{The laws ruling ships are a complex field of legal studies we cannot expand much upon. Most of this is based on ancient custom, and British common law.  Here are some important hooks that you should remember / research. We will deal with some aspects of International Law later.
    }
    \mytext{
    The law of the seas is part common law (tradition), and part conventions (contracts between states).
    \begin{itemize}
        \item concept of international waters: right of passage
        \item ships sail under a flag: law of that country governs ship
        \item international ship register
        \item re-naming or re-painting of ships
        \item positioning systems, navigation (GPS etc.)
        \item AIS automatic identification system (transponders)
        \item convenience flag
        \item false flag, piracy
        \item passage planning, canals, choke points
    \end{itemize}
    }
    
\end{frame}


\begin{frame}{Import Duties}
    \mycomment{
        Goods arriving in ports are often crossing an international border; import duties, such as tariffs, will occur. The declaration (paperwork) will have been included by the seller (unless EXW).
    }
    \showmytikz[1,2,3,4]{../shared/tikz/tikz-freight}
\end{frame}

%%%% ARRIVAL AT PORT / TARIFFS / DOCUMENTS



\begin{frame}{Customs Documents}
\mycomment{
The importer will have to present documents to the local authorities. These will include information about the nature of the imported goods and their value. The authorities will likely check this paperwork, all of which takes time, and are likely to impose fees such as tariffs, excises etc.
}
\mytext{
    \begin{itemize}
        \item Proforma invoice
        \item Commercial invoice
        \item Export packing list
        \item Certificates of origin
        \item Certificate of free sale
        \item Shipper's letter of instruction
        \item Inland bill of landing
        \item Ocean bill of landing
        \item Air waybill
        \item Letter of credit 
        \item ...
    \end{itemize}
}
\end{frame}

\begin{frame}{Tariffs}
\mycomment{
    Tariffs can be levied per unit or per value of the imported goods. 
    
    Within the EU, there is a unified customs area, so there are no tariffs  and no border checks within the EU. 
}
\mytext{
    \begin{itemize}
        \item EU customs area Art. 28 TFEU 
        \item TARIC Tarif intégré des Communautés européennes
        \item $10+4$ digits
    \end{itemize}
}
\end{frame}

\begin{frame}{Excise and VAT}
\mycomment{
    Goods imported may be subject to excises in the importer country, which is a tax on the goods, such as excises on tobacco, alcohol and so forth. These are special taxes levied upon particular goods, while VAT (value added tax) is a general tax on the value of all goods. There may or may not be VAT at the border - a complicated issue. If you have payed VAT in the exporting country, you may be able to subtract it. Otherwise, you will be liable to pay VAT on the sale of the good inside the country.
}
\mytext{
    \begin{itemize}
        \item Excise `Verbrauchssteuer'
        \item EMCS Excise Movement and Control System (EU) \textit{within} EU
        \item Value Added Tax (between companies)
        \item $\neq$ mail order
    \end{itemize}
}
\end{frame}

\begin{frame}{Import Quotas}
\mycomment{
    Sometimes, states limit the number of goods that can be imported to protect their own industry. Then the importer must present the permission to import this particular good. How this quoata is allocated - may per auction? - is a good question.

    States may also ban the import of goods altogether or restrict the import in any other way, and for any kind of reason. Think about sanctions, arms import, drugs, etc.
}
\mytext{
    \begin{itemize}
        \item Cap on absolute number of imports
        \item Allocation?
    \end{itemize}
}
\end{frame}

\begin{frame}{Other Non-tariff Barriers to Trade}
\mycomment{
    On top of that, countries can impose a number of further barriers to trade. Most importantly are standards and regulations: cars may have to have the steering wheel on one side or the other, packaging and documentation may have to be provided in the local language, quality must be proven, etc. etc.
}
\mytext{
    \begin{itemize}
        \item Regulation of all kind
        \item Standards
        \item Composition
        \item Place of Value Added etc.
        \item How to control?
        \item Registration requirements
        \item Certification requirements
        \item ...
    \end{itemize}
}
\end{frame}

\begin{frame}{Classifying Goods}
\mycomment{
    Goods imported must typically be declared, and these declarations follow certain schemes. 
}
\mytext{
    \begin{itemize}
        \item UNSD - ISIC
        \item EUROSTAT - Statistical Classification of Products by Activity – CPA
    \end{itemize}
}
\end{frame}


%%% DUE DILIGENCE


\begin{frame}{Rana Plaza Bangladesh}
\mytext{
        In 2013, a textile factory collapsed. $>1000$ workers died. 
        \showmybitmap{../shared/bitmaps/ranaplaza.png}
}
    \mycomment{
        This sparked outrage through the world and drew fresh attention to the working condition in many factories. International trade separates the consumer from the producer. Many raw materials, but also merchandise, has dubious origins, since the cheapest seller sets the price.

        In Bhopal, India, a factory for pesticides (Methylisocyanat) exploded in 1985 , killing over 10.000 people in the aftermath.
    }
\end{frame}

\begin{frame}{Corporate Due Diligence, CSDDD}
    \mycomment{
    CSDDD is a field much discussed, with much development. How can the consumer and the trader and the producer etc. know and control or influence the working condition and environmental impacts of earlier stages in the value chain? 
    }
    \mytext{
    \begin{itemize}
        \item Lieferkettensorgfaltspflichtengesetz (LkSG) in Germany
        \begin{itemize}
            \item $>1000$ employees (in D)
        \end{itemize}
        \item Corporate Sustainability Due Diligence Directive (CSDDD) EU 
        \item Corporate Sustainability Reporting Directive (CSRD)
    \end{itemize}
    }
\end{frame}

% \begin{frame}{CBAM Carbon Border Adjustment Mechanism}
%     \mycomment{
%     The pricing of \ce{CO2}-emissions creates incentives to move the emission rich production stages abroad.
%     }
%     \mytext{
%         \begin{itemize}
%             \item Transitional phase up to 2026
%             \item cement, iron and steel, aluminium, fertilisers, electricity and hydrogen
%             \item CBAM certificates (national authorities)
%             \item declare yearly GHG (green house gas) emissions
%         \end{itemize}
%     % }
%         % \mycomment{
        
%         Example limestone (cement):
%         - burning for more than 10hrs above 800°C
    
    
%         \ce{CaCO3 -> CaO + CO2}
    
%         \medskip
    
%         Example hydrogen: steam reforming
        
    
%         \ce{CH4 + H2O -> CO + 3H2} (endotherm)
        
%         \ce{CO + H2O -> CO2 + H2} (exotherm)
%     }
% \end{frame}

% \begin{frame}{Carbon Border Emission Trading}
%     \mytext{
    
%     Linked to EU ETS European Union Emissions Trading System

%     \href{https://tradingeconomics.com/commodity/carbon}{Link to Chart}

%     Cap-and-trade system since 2005

%     National allocation plan

%     \medskip

%     Global offsetting 
    
%     CER Certified Emission Reduction, 
    
%     CDM Clean Development Mechanism (UN)
%     }
% \end{frame}

\begin{frame}{Critical Minerals, Conflict Minerals (3TG)}
    \mycomment{
        Some minerals have been linked to the finance of conflicts, in that belligerent parties fight over access to these minerals and use the proceedings from the mines to finance their wars. Working conditions in those mines are often abysmal. The UN has therefore classified some minerals as `conflict minerals'. 
        }
    \mytext{
    \begin{itemize}
        \item tin
        \item tantal
        \item tungsten
        \item gold
    \end{itemize}
    Also: diamonds, ... oil?
    \begin{itemize}
        \item EU Conflict Minerals Regulation \cite{euconflictmineralsregulationexplained}
        \item OECD
        \item private initiatives
    \end{itemize}
    }
\end{frame}


% \begin{frame}{Energy}
% \showmybitmap{../shared/bitmaps/energyimportsnet.png}
% \mycomment{
%     Energy is a good, not a service
%     \begin{itemize}
%         \item Electricity 
%         \item Coal, oil and gas
%         \item But: certificates
%     \end{itemize}
% }
% \end{frame}

% \begin{frame}{Recap: Energy Physics}
% \mytext{
%     \begin{equation}
%         En=Ex+An
%     \end{equation}
%     Energy is exergy (useful) and anergy ('loss')

%     Energy is always preserved

%     Two types of energy, transformable
%     \begin{itemize}
%         \item potential (in a field, also: 'chemical energy')
%         \item kinetic (movement, also: heat)
%     \end{itemize}

%     Transformation of matter into energy per $E=mc^2$ negligible (on earth)
% }
% \end{frame}


% \begin{frame}{Sankey Diagram}

%     \showmytikz{../shared/tikz/tikz-sankey}

%     \href{https://ec.europa.eu/eurostat/cache/sankey/energy/sankey.html?geos=EU27_2020}{https://ec.europa.eu/eurostat/cache/sankey/energy/sankey.html?geos=EU27\_2020}

%     \href{https://www.iea.org/data-and-statistics/data-tools/energy-sankey}{https://www.iea.org/data-and-statistics/data-tools/energy-sankey}
% \end{frame}


% \begin{frame}{Oil Price, Energy Price}
%     \mytable{
%         \begin{tabular}{lr}
%             \toprule
%             Region &   \% \\
%             \midrule
%             USA &  14.7\% \\
%             Saudi Arabia & 13.2\% \\
%             Russia & 12.7\% \\
%             Canada & 5.6\%\\
%             Iraq & 5.5\%\\
%             \bottomrule
%         \end{tabular}
        
%         Crude oil production \textbf{ca. 87.75 million barrel per day}
%     }
% \end{frame}
% \begin{frame}{\ce{CO2} Footprint of Oil per Day}
%      \mytext{
%         \begin{multline}
%              80.75 \cdot 10^{6} bbl \\
%             \approx 12.8 \cdot 10^{9} l \\
%             \approx 12.8 \cdot 10^{6} m^3\\
%             \rightarrow 28.16 \cdot 10^{9} kg\; \ce{CO2}\\
%             \rightarrow 28.16 \cdot 10^{6} t\; \ce{CO2}\\
%         \end{multline}
    
     
%         Source: U.S. Energy Information Administration \href{https://www.eia.gov/international/data/world/petroleum-and-other-liquids/annual-petroleum-and-other-liquids-production}{https://www.eia.gov/international/data/world/petroleum-and-other-liquids/annual-petroleum-and-other-liquids-production}
%      }
% \end{frame}


% \begin{frame}{Daily Oil Output in Tankers}
%     \mytext{
%     A Panamax tanker has typically 60,000 - 80,000 DWT

%     A barrel ($159.987 l$) weights about $131.27 kg$  

%     \bigskip

%     \begin{multline}
%         n_{Panamax, approx.} =  
%         \frac{
%             80.75 \cdot 10^6 bbl \cdot 131kg
%         }{
%             70 \cdot 10^{3}\cdot 10^3
%         }\\
%         \approx 
%         \frac{
%             1146.25 \cdot 10^{6}
%         }{
%             70 \cdot 10^{6}
%         } 
%         \approx
%             151.11\\
%     \end{multline}

%     \vspace{1ex}

%     So the world produces approximately\\
%     151 medium sized oil tankers full \textbf{per day.}

%     }
%     \mycomment{
    
    
%     BP World Energy Outlook (2023) \cite{bpworldenergyoutlook23}

%     \href{https://www.bp.com/content/dam/bp/business-sites/en/global/corporate/pdfs/energy-economics/energy-outlook/bp-energy-outlook-2023.pdf}{https://www.bp.com/content/dam/bp/business-sites/en/global/corporate/pdfs/energy-economics/energy-outlook/bp-energy-outlook-2023.pdf}
%     \href{bp-energy-outlook-2023.pdf}{(LOCAL)}

%     \bigskip

%     \href{https://tradingeconomics.com/commodity/crude-oil}{https://tradingeconomics.com/commodity/crude-oil}
%     }
    
% \end{frame}



% \begin{frame}{Energy Consumption and Growth}
%     \showmybitmap{../shared/bitmaps/energy-vs-gdp-eea.png}
    
%     \mycomment{
    
%     Source: \href{https://www.eea.europa.eu/data-and-maps/figures/correlation-of-per-capita-energy}{https://www.eea.europa.eu/data-and-maps/figures/correlation-of-per-capita-energy}
%     }
% \end{frame}

